\begin{table}[H]
\centering
\caption{Lee (2009) Bounds on ITT Estimates}
\label{tab:lee_bounds}
\begin{threeparttable}
\begin{tabular}[t]{lcc}
\toprule
 & Kenya & Liberia \\
\midrule
Control attrition rate & 0.314 & 0.341 \\
Treatment attrition rate & 0.229 & 0.340 \\
Attrition diff. (T $-$ C) & -0.085 & -0.001 \\
\addlinespace
OLS ITT &   0.031 &  -0.212 \\
\addlinespace
Lee lower bound &  -0.023 &  -0.212 \\
Lee upper bound &   0.154 &  -0.212 \\
\bottomrule
\end{tabular}
\begin{tablenotes}[para,flushleft]
\footnotesize
\item \textit{Notes:} Attrition is defined as having a baseline score but no endline score. Lee (2009) bounds trim the distribution of endline scores in the arm with lower attrition to equalize attrition rates, producing worst-case and best-case ITT estimates. All specifications control for EB ability and strata FE\@. The table reports point-estimate bounds; standard errors are not shown.
\end{tablenotes}
\end{threeparttable}
\end{table}
